% CN 3.15 Exercises (p.110--120)

% ======================================================================
\begin{exercise}{1}
    Recall the Collatz function from p.55 in \S2.9. Prove by cases that, for all
    positive integers $n$,
    \[
        \Collatz^{(2)}(n) \le \tfrac32 n + 1.
    \]
\end{exercise}
\begin{solbox}
\begin{proof}[Proof (by cases)]
    Let $\Collatz(x) = C(x)$.

    \textbf{Case 1: $x$ is odd, $x=2n+1$.}
    \begin{align*}
        C(x) &= 3(2n+1)+1 = 6n+4 = 2(3n+2) \quad \text{(even)} \\
        C^{(2)}(x) &= 3n+2 = 3\left(\frac{x-1}{2}\right)+2
            = \frac{3x-3}{2}+2 = \frac{3}{2}x+\frac12
    \end{align*}
    \[
        \implies C^{(2)}(x) < \frac32 x+1. \tag{1}
    \]

    \textbf{Case 2: $x$ is even, $x=2m$.} Then $C(x)=m$.

    \textbf{Case 2a: $m$ is odd.}
    \[
        C^{(2)}(x) = 3m+1 = \frac32 x+1. \tag{2}
    \]

    \textbf{Case 2b: $m$ is even.}
    \[
        C^{(2)}(x) = \frac{m}{2} = \frac{x}{4}
            < \frac32 x < \frac32 x+1. \tag{3}
    \]

    By (1), (2), (3), in all cases,
    \[
        C^{(2)}(x) \le \frac32 x+1.
    \]
\end{proof}
\end{solbox}

\begin{exercise}{2}
    A set of strings $L$ is \emph{hereditary} if: for every nonempty string $x$
    in $L$, there is a character in $x$ which can be deleted from $x$ to give
    another string in $L$. Prove by contradiction that every nonempty hereditary
    set of strings contains the empty string.
\end{exercise}
\begin{solbox}
\begin{proof}[Proof (by contradiction)]
    Assume that $L$ is a nonempty hereditary set of strings that does
    \emph{not} contain the empty string.

    Let $S = \{\,|x| : x \in L\,\}$ be the set of lengths of strings
    in $L$. Since $L$ is nonempty, $S$ is a nonempty subset of $\N$.
    By the Well-Ordering Principle, $S$ has a smallest element.

    Let $x_0 \in L$ be a string of minimum length, i.e.
    \[
        \forall x \in L,\ |x_0| \le |x|.
    \]
    Since $x_0$ is a nonempty string, its length is
    $|x_0| = n_0 \ge 1$.

    Since $L$ is hereditary and $x_0 \in L$ is nonempty, there is a
    character in $x_0$ which can be deleted to give another string
    $x_1 \in L$, where
    \[
        |x_1| = n_0 - 1.
    \]

    Since $x_0$ is the smallest string in $L$, we have
    $|x_0| \le |x|$ for every $x \in L$. But
    \[
        |x_1| = n_0 - 1 < n_0 = |x_0|,
    \]
    which contradicts the minimality of $x_0$.

    This is a contradiction. Therefore, the assumption is false, and
    $L$ must contain the empty string.
\end{proof}
\end{solbox}

\begin{exercise}{3}
    Review the statements of all the theorems (including corollaries) in the
    previous chapters.
    \begin{enumerate}[label=(\alph*)]
        \item Which ones are implications?
        \item Which ones are equivalences?
        \item Which ones are existential statements?
        \item Which ones are universal statements?
    \end{enumerate}
\end{exercise}

\begin{exercise}{4}
    Review the proofs of all the theorems in the previous chapters.
    \begin{enumerate}[label=(\alph*)]
        \item Which ones use proof by construction?
        \item Which ones prove an implication by proving its contrapositive?
        \item Which ones use proof by cases?
    \end{enumerate}
\end{exercise}

% ======================================================================
\begin{exercise}{5}
    Let $E$ be a finite set and $\I$ a collection of subsets of $E$ satisfying:
    \begin{itemize}
        \item[(I1)] $\emptyset\in\I$.
        \item[(I2)] (hereditary) Whenever $X\subseteq Y$ and $Y\in\I$, we also
                    have $X\in\I$.
        \item[(I3)] (independence augmentation) Whenever $X,Y\in\I$ and
                    $|X|<|Y|$, there exists $y\in Y\setminus X$ such that
                    $X\cup\{y\}\in\I$.
    \end{itemize}
    The sets in $\I$ are called \emph{independent}. E.g.\ for $E=\{a,b,c\}$, each
    of the following satisfies (I1)--(I3):
    $\{\emptyset,\{a\},\{b\},\{a,b\}\}$;
    $\{\emptyset,\{a\},\{b\},\{c\}\}$;
    $\{\emptyset,\{a\},\{b\},\{c\},\{a,c\},\{b,c\}\}$;
    $\{\emptyset,\{a\},\{b\},\{c\},\{a,b\},\{a,c\},\{b,c\}\}$;
    $\{\emptyset\}$.
    Counterexamples:
    $\{\{a\},\{b\},\{a,b\}\}$ fails (I1), (I2);
    $\{\emptyset,\{a\},\{a,b\}\}$ fails (I2);
    $\{\emptyset,\{a\},\{b\},\{c\},\{a,b\}\}$ fails (I3);
    $\emptyset$ fails (I1).
    \begin{enumerate}[label=(\alph*)]
        \item For each counterexample above that fails (I2) and/or (I3), give
              specific sets $X$ and $Y$ for which the property fails.
        \item Prove that, for a social network, the set of cliques satisfies (I1)
              and (I2) but not, in general, (I3).
        \item Prove by contradiction that all maximal independent sets have the
              same size.
    \end{enumerate}
    Now give each $e\in E$ a weight $w(e)\ge0$, and define
    $w(X):=\sum_{e\in X}w(e)$. We seek an independent set of maximum weight
    using the \emph{greedy algorithm}:
    \begin{quote}
        $X:=\emptyset$, $i:=1$.\\
        Repeat: if there exists $y\in E\setminus X$ with $X\cup\{y\}\in\I$, choose
        such a $y$ with largest weight, call it $e_i$, add it to $X$ and increase
        $i$; else stop and output $X$.
    \end{quote}
    \begin{enumerate}[label=(\alph*),start=4]
        \item Construct a simple social network with nonnegative weights on the
              people such that the greedy algorithm, with $\I$ being the set of
              all cliques, does not find the maximum weight clique.
    \end{enumerate}
    Now suppose $\I$ satisfies (I1)--(I3), and the greedy algorithm chooses
    $e_1,e_2,\dots,e_r$ in that order.
    \begin{enumerate}[label=(\alph*),start=5]
        \item Prove by contradiction that
              $w(e_1)\ge w(e_2)\ge\cdots\ge w(e_r)$.
        \item Prove that the greedy algorithm finds a maximum weight independent
              set.
    \end{enumerate}
\end{exercise}
\begin{solbox}[long]

\begin{solution}[(a)]
    % TODO: give specific $X,Y$ for each counterexample that fails
    % (I2) and/or (I3).
\end{solution}

\begin{solution}[(b)]
\begin{proof}[Proof of (b)]
    Let $E$ = vertices of a social network, and let $\I$ = the set of
    all cliques.

    \medskip
    \textbf{(I1).} Let $X$ be a clique, and consider $X = \emptyset$.
    Since $\emptyset$ contains no elements, there do not exist
    $u,v\in\emptyset$ with $u\ne v$.
    \[
        \text{So the implication ``}u\ne v \implies (u,v)\text{ is
        edge''}
    \]
    holds vacuously for all $u,v\in\emptyset$. So $\emptyset$ is a
    clique, i.e.
    \[
        \emptyset\in\I.
    \]

    \medskip
    \textbf{(I2).} Need to prove: if $X\subseteq Y$ and $Y$ is a
    clique, then $X$ is a clique.

    For all $u,v\in X$ with $u\ne v$: since $X\subseteq Y$,
    \[
        u,v\in Y \quad \text{with } u\ne v.
    \]
    Since $Y$ is a clique, $(u,v)$ is an edge. Hence $X$ is also a
    clique.

    \medskip
    \textbf{(I3) fails in general.} Consider the graph with vertices
    $\{a,b,c,d\}$ and a single edge $(b,c)$. Let
    \[
        X=\{a\}, \qquad Y=\{b,c\}.
    \]
    $X$ and $Y$ are cliques, and $Y\setminus X = \{b,c\}$. Checking
    each candidate:
    \[
        X\cup\{b\} = \{a,b\} \text{ is not a clique, since } (a,b)
        \text{ is not an edge.}
    \]
    \[
        X\cup\{c\} = \{a,c\} \text{ is not a clique, since } (a,c)
        \text{ is not an edge.}
    \]
    So no valid $y\in Y\setminus X$ exists, and (I3) fails.
\end{proof}

\begin{remark}
    This is the key place where the ``clique matroid'' fails to be a
    genuine matroid: in the vector-space analogy, (I3) is the
    exchange axiom that guarantees any linearly independent set can
    be extended toward a larger independent set. Cliques do not have
    this exchange property in general, which is exactly why the
    greedy algorithm can fail on them (part (d)).
\end{remark}
\end{solution}

\begin{solution}[(c)]
\begin{proof}[Proof of (c)]
    An independent set $X\in\I$ is \emph{maximal} if
    \[
        \nexists\, Y\in\I \text{ such that } X\subsetneq Y.
    \]

    We need to prove:
    \[
        \forall X,Y\in\I,\quad X \text{ maximal} \wedge Y \text{
        maximal} \implies |X|=|Y|.
    \]

    Assume, for contradiction, there exist two maximal independent
    sets $X,Y\in\I$ with $|X|\ne|Y|$. WLOG, assume
    \[
        |X|<|Y|.
    \]

    Since $X,Y\in\I$ and $|X|<|Y|$, by (I3) there exists
    $y\in Y\setminus X$ such that
    \[
        X\cup\{y\}\in\I.
    \]

    Since $y\notin X$, we have
    \[
        X\subsetneq X\cup\{y\}, \qquad X\cup\{y\}\in\I.
    \]
    This contradicts the assumption that $X$ is maximal.

    Therefore, the assumption is false: all maximal independent sets
    have the same size.
\end{proof}

\begin{remark}
    This is the exact analogue of the fact that any two bases of a
    vector space have the same size (i.e.\ dimension is well
    defined). There, the argument uses the Steinitz exchange lemma;
    here (I3) plays precisely that role.
\end{remark}
\end{solution}

\begin{solution}[(d)]
    % TODO: construct a concrete social network graph and weights
    % where greedy fails to find the maximum weight clique.
\end{solution}

\begin{solution}[(e)]
\begin{proof}[Proof of (e)]
    When $\I$ satisfies (I1)--(I3), the greedy algorithm chooses
    elements of $E$ in a specific order $e_1,e_2,\dots,e_r$. We prove
    \[
        w(e_i) \ge w(e_{i+1}) \qquad \text{for all } 1\le i<r.
    \]

    Recall $e_i$ is the maximum-weight element among all
    $y\in E\setminus X_{i-1}$ satisfying $X_{i-1}\cup\{y\}\in\I$:
    \[
        X_{i-1}\cup\{e_i\}\in\I,
    \]
    \[
        \forall y \text{ with } X_{i-1}\cup\{y\}\in\I,\quad
        w(e_i)\ge w(y).
    \]

    Assume, for contradiction, that
    \[
        w(e_i) < w(e_{i+1}) \qquad \text{for some } 1\le i<r.
    \]

    Since $e_{i+1}$ was chosen by the algorithm at round $i+1$,
    \[
        X_i \cup \{e_{i+1}\} \in \I.
    \]
    Since $X_{i-1}\subseteq X_i$,
    \[
        X_{i-1}\cup\{e_{i+1}\} \subseteq X_i\cup\{e_{i+1}\}.
    \]
    Since $X_i\cup\{e_{i+1}\}\in\I$, by (I2),
    \[
        X_{i-1}\cup\{e_{i+1}\}\in\I.
    \]

    This means $e_{i+1}$ is one of the candidates $y$ satisfying
    $X_{i-1}\cup\{y\}\in\I$ at round $i$. Since $e_i$ was chosen at
    round $i$ as the maximum-weight element among all such
    candidates,
    \[
        w(e_i) \ge w(e_{i+1}),
    \]
    which contradicts the assumption $w(e_i) < w(e_{i+1})$.

    Therefore the assumption is false, and
    \[
        w(e_i)\ge w(e_{i+1}) \qquad \text{for all } 1\le i<r.
    \]
\end{proof}

\begin{remark}
    This mirrors a standard fact from linear algebra: if you build up
    an independent set greedily by always adding the ``largest''
    remaining valid candidate, the sequence of choices is
    automatically monotone --- the exchange property (I2) is what
    lets you ``pull'' a later candidate back to an earlier round
    without breaking independence.
\end{remark}
\end{solution}

\begin{solution}[(f)]
\begin{proof}[Proof of (f)]
    Let $Y$ be a maximum weight independent set. By the stopping
    condition of the greedy algorithm, when $X=\{e_1,\dots,e_r\}$,
    there is no $y\in E\setminus X$ such that
    \[
        X\cup\{y\}\in\I.
    \]
    So $X$ is a maximal independent set.

    By part (c), we know
    \[
        |X|=|Y|=r.
    \]
    Write $Y=\{f_1,\dots,f_r\}$ with
    \[
        w(f_1)\ge w(f_2) \ge \cdots \ge w(f_r).
    \]

    \medskip
    \textbf{Main claim.} For each $i=1,\dots,r$,
    \[
        w(e_i)\ge w(f_i).
    \]

    By contradiction (minimal counterexample): assume there exists a
    minimum $i$ such that $w(e_i) < w(f_i)$. By minimality of $i$,
    \[
        w(e_k)\ge w(f_k) \qquad \text{for all } k<i.
    \]

    Let
    \[
        X_{i-1}=\{e_1,\dots,e_{i-1}\}, \qquad F_i=\{f_1,\dots,f_i\}.
    \]
    Since $X_{i-1}\subseteq X\in\I$, by (I2), $X_{i-1}\in\I$. Since
    $F_i\subseteq Y\in\I$, by (I2), $F_i\in\I$. Also
    \[
        |X_{i-1}|=i-1 < i = |F_i|.
    \]

    By (I3), there exists $y\in F_i\setminus X_{i-1}$ such that
    \[
        X_{i-1}\cup\{y\}\in\I.
    \]

    Since $y\in F_i = \{f_1,\dots,f_i\}$, and these are sorted in
    decreasing order of weight, $f_i$ has the smallest weight among
    them, so
    \[
        w(y) \ge w(f_i).
    \]

    Now, since $y\notin X_{i-1}$ and $X_{i-1}\cup\{y\}\in\I$, $y$ is a
    valid candidate at round $i$ of the greedy algorithm. But $e_i$
    was chosen at round $i$ as the element of maximum weight among
    all such valid candidates, so
    \[
        w(e_i) \ge w(y) \ge w(f_i).
    \]
    This contradicts the assumption $w(e_i) < w(f_i)$. Hence no such
    $i$ exists, and
    \[
        w(e_i)\ge w(f_i) \qquad \text{for all } i=1,\dots,r.
    \]

    \medskip
    Summing over $i$:
    \[
        w(X) = \sum_{i=1}^r w(e_i) \ge \sum_{i=1}^r w(f_i) = w(Y).
    \]
    Since $Y$ was assumed to be a maximum weight independent set,
    $w(X)\ge w(Y)$ forces
    \[
        w(X)=w(Y).
    \]
    Therefore $X$ is also a maximum weight independent set.
    \[
        \therefore \text{Greedy algorithm finds a maximum weight
        independent set.}
    \]
\end{proof}

\begin{remark}
    This whole argument is the discrete, combinatorial analogue of
    the Matroid Greedy Theorem, itself a generalisation of the fact
    that Gaussian elimination / the Gram--Schmidt-style greedy basis
    selection process finds an optimal basis with respect to any
    linear weighting --- the same ``minimal counterexample +
    exchange'' technique (Steinitz exchange lemma) underlies both.
\end{remark}
\end{solution}

\FloatBarrier
\end{solbox}

\begin{exercise}{6}
    Prove by mathematical induction: for all $k$,
    \[
        (*)\qquad \text{the sum of the first $k$ odd numbers equals } k^2.
    \]
    \begin{enumerate}[label=(\alph*)]
        \item Give a simple expression for the $k$-th odd number.
        \item Inductive basis: prove $(*)$ for $k=1$.
        \item Inductive step: let $k\ge1$ and assume $(*)$ for $k$. Express the
              sum of the first $k+1$ odd numbers in terms of the sum of the first
              $k$ odd numbers, plus something else.
        \item Use the inductive hypothesis to replace the sum of the first $k$ odd
              numbers by something else.
        \item Simplify your expression. What do you notice?
        \item Draw your final conclusion, stating that you are using the
              Principle of Mathematical Induction.
    \end{enumerate}
\end{exercise}

\begin{exercise}{7}
    Prove by induction on $n$ that, for all $n\in\N_0$,
    \[
        1+2+2^2+2^3+\cdots+2^n = 2^{n+1}-1.
    \]
\end{exercise}

\begin{exercise}{8}
    Prove by induction that, for all $n$:
    \[
        1^2+2^2+\cdots+n^2 = \frac{n(n+1)(2n+1)}{6}.
    \]
\end{exercise}

\begin{exercise}{9}
    Prove by induction that, for all $n$:
    \[
        1^3+2^3+\cdots+n^3 = (1+2+\cdots+n)^2.
    \]
\end{exercise}

% ======================================================================
\begin{exercise}{10}
    Consider the following algorithm:
    \begin{quote}
        for each $i=1,2,\dots,n$:\\
        \hspace*{1em} for each $j=1,2,\dots,i$:\\
        \hspace*{2em} for each $k=1,2,\dots,j$:\\
        \hspace*{3em} [some action]
    \end{quote}
    \begin{enumerate}[label=(\alph*)]
        \item For each pair $i,j$, how many times is the action performed?
        \item For each $i$, how many times is the action performed? Write this
              both as a sum and (using a theorem in this chapter) as a simple
              algebraic expression in $i$.
        \item Write a sum, in terms of $n$, for the total number of times the
              action is performed.
        \item Work out this expression for $n=1,2,3,4,5$ (and more).
        \item Explore how these numbers behave.
        \item Conjecture a simple algebraic expression in $n$ for the total.
        \item Prove, by induction on $n$, that your expression is correct.
    \end{enumerate}
\end{exercise}
\begin{solbox}[long]
\begin{solution}
    \begin{enumerate}[label=(\alph*)]
    \item For fixed $i,j$: $j$ loop from $1$ to $i$, $k$ loop from $1$ to
        $j$. Observe that for each pair $(i,j)$, the innermost action is
        performed $j$ times.

    \item For fixed $i$: $j$ loops from $1$ to $i$, and for each $j$ the
        inner loop runs $j$ times. So the number of times the action is
        performed is
        \[
            \sum_{j=1}^{i} j = \frac{i(i+1)}{2}.
        \]

    \item Summing over $i$ from $1$ to $n$:
        \[
            \text{num of action} = \sum_{i=1}^{n}\sum_{j=1}^{i} j
            = \sum_{i=1}^{n} \frac{i(i+1)}{2}.
        \]

    \item Evaluating $\displaystyle\sum_{i=1}^{n}\frac{i(i+1)}{2}$ for
        small $n$:
        \[
            n=1: 1, \quad n=2: 4, \quad n=3: 10, \quad n=4: 20, \quad
            n=5: 35.
        \]

    \item The sequence $1,4,10,20,35,\dots$ is the sequence of
        \emph{tetrahedral numbers}. Each term is obtained from the
        previous by adding the next triangular number
        ($1,3,6,10,15,\dots$), and the values grow roughly cubically in
        $n$.

    \item Conjecture:
        \[
            \sum_{i=1}^{n}\frac{i(i+1)}{2} = \frac{n(n+1)(n+2)}{6}.
        \]
        \emph{Derivation:} using $\displaystyle\sum_{i=1}^n i^2 =
        \frac{n(n+1)(2n+1)}{6}$ and $\displaystyle\sum_{i=1}^n i =
        \frac{n(n+1)}{2}$,
        \begin{align*}
            \sum_{i=1}^n \frac{i^2+i}{2}
            &= \frac12\left[\frac{n(n+1)(2n+1)}{6}
                + \frac{n(n+1)}{2}\right] \\
            &= \frac{1}{12}\left[n(n+1)(2n+1) + 3n(n+1)\right] \\
            &= \frac{1}{12}\,n(n+1)\big[(2n+1)+3\big]
                = \frac{1}{12}\,n(n+1)(2n+4) \\
            &= \frac{1}{6}\,n(n+1)(n+2).
        \end{align*}
        This matches the numerical values in part (d): $n=5$ gives
        $\frac{5\cdot6\cdot7}{6}=35$.\ \checkmark

    \item \textbf{Proof (by induction on $n$).}

        \emph{Base case ($n=1$).} LHS $=\dfrac{1\cdot2}{2}=1$.
        RHS $=\dfrac{1\cdot2\cdot3}{6}=1$. Equal, so the base case holds.

        \emph{Inductive hypothesis.} Assume for some $n\ge1$,
        \[
            \sum_{i=1}^{n}\frac{i(i+1)}{2} = \frac{n(n+1)(n+2)}{6}.
        \]

        \emph{Inductive step.} We show the formula holds for $n+1$:
        \begin{align*}
            \sum_{i=1}^{n+1}\frac{i(i+1)}{2}
            &= \sum_{i=1}^{n}\frac{i(i+1)}{2} + \frac{(n+1)(n+2)}{2} \\
            &= \frac{n(n+1)(n+2)}{6} + \frac{(n+1)(n+2)}{2}
                \qquad \text{(by inductive hypothesis)} \\
            &= \frac{n(n+1)(n+2) + 3(n+1)(n+2)}{6} \\
            &= \frac{(n+1)(n+2)(n+3)}{6}.
        \end{align*}
        This matches the conjectured formula with $n$ replaced by $n+1$.

        By the Principle of Mathematical Induction, the formula
        \[
            \sum_{i=1}^{n}\frac{i(i+1)}{2} = \frac{n(n+1)(n+2)}{6}
        \]
        holds for all $n\ge1$. \qed
    \end{enumerate}
\end{solution}
\end{solbox}

\begin{exercise}{11}
    The $n$-th harmonic number is $H_n = 1+\frac12+\frac13+\cdots+\frac1n$. We
    prove by induction that $H_n\ge\log_e(n+1)$.
    \begin{enumerate}[label=(\roman*)]
        \item Inductive basis: prove that $H_1\ge\log_e(1+1)$.
        \item Let $n\ge1$ and assume $H_n\ge\log_e(n+1)$. Write $H_{n+1}$
              recursively using $H_n$, use the inductive hypothesis, then use
              $\log_e(1+x)\le x$ and an elementary property of logarithms to show
              $H_{n+1}\ge\log_e(n+2)$.
        \item What can you now conclude, and why?
    \end{enumerate}
    Afterthoughts: try to prove $H_n\ge\log_e n$ directly by induction --- where
    is the snag? Would the proof work for other bases? Can you prove
    $H_n\le(\log_e n)+1$?
\end{exercise}
\begin{solbox}[long]
\begin{solution}
    The $n$-th harmonic number $H_n$ is defined by
\[
    H_n = 1+\frac12+\frac13+\cdots+\frac1n.
\]

\section*{Main result: $H_n \ge \ln(n+1)$}

\begin{proof}[Proof (by induction on $n$)]
    \textbf{Base case ($n=1$).} $H_1=1$, and $\ln(1+1)=\ln2<1$. So
    $H_1 \ge \ln(1+1)$.

    \textbf{Inductive hypothesis.} Assume $H_n \ge \ln(n+1)$ for some
    $n\ge1$.

    \textbf{Inductive step.} We show $H_{n+1}\ge\ln(n+2)$.

    By definition,
    \[
        H_{n+1} = H_n+\frac{1}{n+1} \ge \ln(n+1)+\frac{1}{n+1}.
    \]
    We claim $\ln(n+1)+\dfrac{1}{n+1} \ge \ln(n+2)$. Observe that
    \[
        \ln(n+2)-\ln(n+1) = \ln\left(\frac{n+2}{n+1}\right)
        = \ln\left(1+\frac{1}{n+1}\right) \le \frac{1}{n+1},
    \]
    using $\ln(1+x)\le x$ (with $x=\frac{1}{n+1}$). Rearranging,
    \[
        \ln(n+1)+\frac{1}{n+1} \ge \ln(n+2).
    \]
    Combining the two inequalities,
    \[
        H_{n+1} \ge \ln(n+1)+\frac{1}{n+1} \ge \ln(n+2).
    \]
    So the statement holds for $n+1$.

    \textbf{Conclusion.} By the Principle of Mathematical Induction,
    $H_n \ge \ln(n+1)$ for all $n\ge1$.
\end{proof}

\section*{Advanced afterthoughts}

\begin{enumerate}
\item The inequality $H_n\ge\ln(n+1)$ implies $H_n\ge\ln n$, since
    $\ln(n+1)\ge\ln n$. Trying to prove $H_n\ge\ln n$ \emph{directly}
    by induction runs into a snag: in the inductive step, one would
    need
    \[
        \ln(n+1)-\ln n = \ln\left(1+\frac1n\right) \le \frac{1}{n+1},
    \]
    but $\ln(1+x)\le x$ only gives the weaker bound
    $\ln(1+\frac1n)\le\frac1n$, which is not tight enough (since
    $\frac1n > \frac{1}{n+1}$). This illustrates that induction
    sometimes requires proving something \emph{stronger} than what
    one set out to prove --- the extra ``room'' in $\ln(n+1)$
    compared to $\ln n$ is exactly what makes the induction go
    through.

\item The proof relies on the inequality $\ln(1+x)\le x$, which is a
    property specific to the natural logarithm (it comes from
    $\ln(1+x)$ being concave with tangent line $y=x$ at $x=0$). For a
    logarithm to another base $b\ne e$, the analogous inequality
    $\log_b(1+x)\le x$ does not hold in general, so the proof does
    not transfer directly to other bases without first converting
    back to $\ln$ via the change-of-base formula.

\item \textbf{Claim: $H_n \le \ln n+1$.}
    \begin{proof}[Proof (by induction on $n$)]
        \textbf{Base case ($n=1$).} $H_1=1$. RHS $=\ln1+1=0+1=1$.
        So $H_1\le\ln1+1$.

        \textbf{Inductive hypothesis.} Assume $H_n\le\ln n+1$ for
        some $n\ge1$.

        \textbf{Inductive step.} We show $H_{n+1}\le\ln(n+1)+1$.

        By definition, $H_{n+1}=H_n+\dfrac{1}{n+1}$. By the
        inductive hypothesis,
        \[
            H_{n+1} \le \ln n+1+\frac{1}{n+1}.
        \]
        It remains to show $\ln n+\dfrac{1}{n+1}\le\ln(n+1)$, i.e.
        \[
            \ln(n+1)-\ln n \ge \frac{1}{n+1}.
        \]
        Using the inequality $\ln x \ge 1-\dfrac1x$ for $x>0$, take
        $x=\dfrac{n+1}{n}$:
        \[
            \ln\left(\frac{n+1}{n}\right) \ge 1-\frac{n}{n+1}
            = \frac{1}{n+1}.
        \]
        Since $\ln\left(\dfrac{n+1}{n}\right)=\ln(n+1)-\ln n$, this
        gives
        \[
            \ln(n+1)-\ln n \ge \frac{1}{n+1}.
        \]
        Combining:
        \[
            H_{n+1} \le \ln n+1+\frac{1}{n+1} \le \ln(n+1)+1.
        \]
        So the statement holds for $n+1$.

        \textbf{Conclusion.} By the Principle of Mathematical
        Induction, $H_n \le \ln n+1$ for all $n\ge1$.
    \end{proof}
\end{enumerate}
\end{solution}
\end{solbox}

\begin{exercise}{12}
    (Challenge) Let $E$ be a finite set and $\mathcal S$ a set of subsets of $E$.
    \begin{itemize}
        \item $\mathcal S$ is \emph{linear under $\triangle$} if
              $X\triangle Y\in\mathcal S$ for all $X,Y\in\mathcal S$.
        \item $\mathrm{span}(\mathcal S) := \{X_1\triangle\cdots\triangle X_k :
              X_1,\dots,X_k\in\mathcal S,\ k\in\N\}$ (the symmetric difference of
              zero sets is $\emptyset$).
        \item $\mathcal S$ is \emph{dependent under $\triangle$} if for some
              distinct $X_1,\dots,X_k\in\mathcal S$,
              $X_1\triangle\cdots\triangle X_k=\emptyset$; otherwise
              \emph{independent}.
        \item A \emph{circuit} is a minimal dependent set; a \emph{base} is a
              maximal independent set.
    \end{itemize}
    \begin{enumerate}[label=(\alph*)]
        \item Let $\mathcal S$ be linear under $\triangle$. Prove by induction on
              $k$ that for all $X_1,\dots,X_k\in\mathcal S$,
              $X_1\triangle\cdots\triangle X_k\in\mathcal S$.
        \item Prove that $\mathcal S$ is linear under $\triangle$ iff
              $\mathcal S=\mathrm{span}(\mathcal S)$.
        \item Prove by induction on $k$ that for all $X_1,\dots,X_k\in\mathcal
              S$, either $X_1\triangle\cdots\triangle X_k=\emptyset$, or there
              exist distinct $Y_1,\dots,Y_n\in\mathcal S$ with $n\le k$ and
              $X_1\triangle\cdots\triangle X_k=Y_1\triangle\cdots\triangle Y_n$.
        \item Prove that $\mathcal S$ is dependent iff there exists
              $X\in\mathcal S$ with $X\in\mathrm{span}(\mathcal S\setminus\{X\})$.
        \item Prove that if $\mathcal S$ is minimal dependent, then for every
              $X\in\mathcal S$, $X\in\mathrm{span}(\mathcal S\setminus\{X\})$.
        \item Prove that every minimal spanning set is independent.
        \item Prove that every minimal spanning set is a maximal independent set.
        \item Prove that every maximal independent set is a spanning set.
        \item Prove that every maximal independent set is a minimal spanning set.
        \item Prove that $\mathcal S$ is a minimal spanning set iff it is a
              maximal independent set.
        \item Prove that if $\mathcal S$ is independent then
              $|\mathrm{span}(\mathcal S)|=2^{|\mathcal S|}$.
        \item Prove the independence augmentation property: if $\mathcal S$ and
              $\mathcal T$ are independent and $|\mathcal T|=|\mathcal S|+1$,
              then there exists $X\in\mathcal T\setminus\mathcal S$ such that
              $\mathcal S\cup\{X\}$ is independent.
        \item Prove that all bases have the same size.
        \item Do all minimal dependent sets have the same size? Prove your claim.
        \item Prove that for every base $\mathcal B$ of $\mathcal S$ and every
              $X\in\mathcal S$, there is a unique way to write $X$ as a symmetric
              difference of members of $\mathcal B$, each used at most once.
        \item For each of (a)--(o), state what type of proof you used.
        \item Design a simple algorithm: INPUT a collection $\mathcal S$ of
              subsets of $E$ with nonnegative weights; OUTPUT an independent
              $\mathcal X\subseteq\mathcal S$ of maximum total weight. Does it
              always find the best solution? Justify using a previous exercise.
    \end{enumerate}
\end{exercise}
\begin{solbox}
\begin{solution}[(a)]

\end{solution}
\begin{solution}[(b)]

\end{solution}
\begin{solution}[(c)]

\end{solution}
\begin{solution}[(d)]

\end{solution}
\begin{solution}[(e)]

\end{solution}
\begin{solution}[(f)]

\end{solution}
\begin{solution}[(g)]

\end{solution}
\begin{solution}[(h)]

\end{solution}
\begin{solution}[(i)]

\end{solution}
\begin{solution}[(j)]

\end{solution}
\begin{solution}[(k)]

\end{solution}
\begin{solution}[(l)]

\end{solution}
\begin{solution}[(m)]

\end{solution}
\begin{solution}[(n)]

\end{solution}
\begin{solution}[(o)]

\end{solution}
\begin{solution}[Algorithm \& justification]

\end{solution}
\end{solbox}

\begin{exercise}{13}
    Identify the errors in the following ``theorem'' and ``proof''.

    \textbf{``Theorem.''} For all $n\in\N$, a circular disc can be cut into $2^n$
    pieces by $n$ straight lines.

    \textbf{``Proof.''}
    \begin{enumerate}
        \item For $n=0$ the disc is in one piece $=2^0$.
        \item For $n=1$, a single line cuts it into two pieces $=2^1$.
        \item For $n=2$, a second line crossing the first cuts across both pieces,
              giving $2\times2=2^2$ pieces.
        \item It is clear from the cases so far that, once the circle is cut by
              some number of lines, another line can be used to divide every
              piece into two, doubling the number of pieces.
        \item So with $n$ lines, repeated doubling gives $2^n$ pieces. ``$\square$''
    \end{enumerate}
\end{exercise}
\begin{solbox}
\begin{solution}

\end{solution}
\end{solbox}

\begin{exercise}{14}
    Identify the errors in the following ``theorem'' and ``proof''.

    \textbf{``Theorem.''} Every rational number is an integer.

    \textbf{``Proof.''}
    \begin{enumerate}
        \item We need to show that the sets $\Z$ and $\Q$ are equal.
        \item To prove two sets equal, we can prove the appropriate subset and
              superset relations.
        \item We first prove $\Z\subseteq\Q$.
        \item Let $x\in\Z$.
        \item Since $x$ is an integer, $x=\frac x1$.
        \item So $x$ is a quotient of integers.
        \item Therefore $x\in\Q$.
        \item So $\Z\subseteq\Q$.
        \item Now we need to prove the superset relation, the reverse of subset.
        \item So we will prove that $\Q\supseteq\Z$.
        \item Let $x\notin\Q$.
        \item Then $x$ cannot be written as $\frac pq$ with $p,q$ integers.
        \item So it certainly cannot be written in this form with $q=1$.
        \item Therefore $x\notin\Z$.
        \item So $\Q\supseteq\Z$.
        \item Hence $\Z=\Q$. ``$\square$''
    \end{enumerate}
\end{exercise}
\begin{solbox}
\begin{solution}

\end{solution}
\end{solbox}

\begin{exercise}{15}
    Let $G$ be the cryptosystem obtained from Caesar slide (Exercise 2.15) by
    restricting the keyspace to $K_G=\{\mathtt a,\mathtt b,\dots,\mathtt m\}$.
    Prove that $G$ is not idempotent.
\end{exercise}
\begin{solbox}
\begin{solution}

\end{solution}
\end{solbox}

